(1) 一个高低折射率周期的光学厚度为 \(\Lambda_{\rm opt}=n_Hd_H+n_Ld_L\)。在法向入射时从相邻等效周期界面返回的两束光多经历一个完整周期的往返光程 \(2\Lambda_{\rm opt}\)，由此可建立反射相位相干增强条件。
(2) 相邻周期反射相位差为 \(\Delta\phi=2\pi(2\Lambda_{\rm opt})/\lambda_0\)，第一阶Bragg共振取 \(\Delta\phi=2\pi\)，所以中心波长 \(\lambda_B=2(n_Hd_H+n_Ld_L)\)。实际精确反射谱的中心与宽度还会取决于各界面反射相位和强度，特别是大折射率对比。
(3) 四分之一波设计要求每种层的单程光学厚度为 \(\lambda_B/4\)，即 \(n_Hd_H=n_Ld_L=\lambda_B/4\)。因此物理厚度 \(d_H=\lambda_B/(4n_H)\)、\(d_L=\lambda_B/(4n_L)\)，代入上式得一个周期总光学厚度 \(\lambda_B/2\)，与Bragg条件一致。
(4) 在设计波长处多个周期的反射振幅相干叠加，增加周期数通常使反射率提高、形成更显著的禁带，但具体反射率必须用光学传输矩阵求出。若 \(n_H\to n_L\)，每个界面的Fresnel反射振幅趋零，即使各周期相位匹配也无可累积反射，禁带消失。
